\documentclass[11pt]{article}
\usepackage[letterpaper,margin=1in]{geometry}
\usepackage{amsmath}
\usepackage[colorlinks=true,urlcolor=blue]{hyperref}
\setlength{\parindent}{0pt}
\setlength{\parskip}{0.6em}
\begin{document}
\section*{Significance Statement}
A correlated insulator has no low-energy electron poles, but its Green function can still vanish along a dispersing curve. We show that this hidden dispersion identifies why the insulator forms. In the Mott limit, its scale directly measures short-range spin correlations; in the charge-ordered limit, local charge-defect counting fixes a different universal coefficient. With both interactions present, the dispersion strengthens toward $V/U=1/4$ and switches branches. Dense momentum sampling independently reveals a simultaneous jump in double occupancy and checkerboard order, collapse of antiferromagnetic correlations, and discontinuity of the ground-state energy derivative. Green-function zeros therefore expose a first-order competition between insulating mechanisms while retaining momentum information after conventional quasiparticle poles have disappeared.
\end{document}
